\section{问题重述}

% 写作提示：用自己的话交代真实背景、已知条件与数据、待求目标、约束和输出；对含糊概念作必要澄清。
% 不要整段复制题面，也不要在问题重述中提前给出模型公式和求解结果。下方内容仅为结构示例，须按赛题改写。

\subsection{问题背景}
情感识别是人工智能、自然人机交互和社会计算的关键支撑技术之一。在实际应用中，通常综合文本语义、语音声学特征和面部表情等多模态信息进行情感判断。然而，复杂场景下的多模态情感识别面临诸多挑战：时序对齐困难、模态缺失普遍、预测结果缺乏可解释性。本文基于CMU-MOSEI英文多模态情感数据集，围绕特征提取、鲁棒预测和可解释性展开系统建模研究。

\subsection{问题内容}
\textbf{问题一：多模态特征提取与时序对齐}

针对附件1提供的100条原始视频样本，建立多模态特征提取与时序对齐的数学模型。需要从文本、语音、视觉三种模态中提取有效特征，实现不同模态时序信息的统一表示，并生成规范的特征文件以支持后续建模。文本模态采用BERT编码（768维），视觉模态包含面部动作单元等特征（35维），音频模态包含COVAREP声学特征（74维）。时序长度在不同样本间差异显著，需要设计对齐策略将可变长度序列映射为固定维度特征向量。

\textbf{问题二：模态缺失下的鲁棒情感预测}

针对局部时刻模态信息缺失的复杂场景，构建鲁棒的情感预测模型。模态局部缺失指文本、语音或视觉中的部分时间片段不可用（全零标记）。模型需要在局部模态信息不完整的情况下保持稳定的情感预测，并对附件3中30个缺失样本进行预测。情感标签为回归值，范围[-3, 3]，负值表示消极，正值表示积极。

\textbf{问题三：可解释性情感预测模型}

构建可解释的多模态情感预测模型。模型需要在给出情感预测的同时，能够定位关键时刻、识别重要模态，并回溯到原始文本片段、音频时刻或视觉关键帧。对附件4中20个专项样本进行可解释性分析，量化各模态对预测结果的贡献度，并可视化关键证据的时空分布。

% 技术路线图、研究框架图和算法流程图一律用 paper-diagram Skill 绘制——
% 优先套用其内置模板，产出可编辑 .drawio + PNG，
% 并用 check_layout.py / export_figure.py 校验后插入正文。

